\begin{abstract}

    Die zunehmende Urbanisierung und Motorisierung führten in den vergangenen Jahren zu einer starken Überlastung städtischer Verkehrsinfrastrukturen, wobei die Parkplatzsuche eine zentrale Rolle spielte. In diesem Kontext gewannen Smart-Parking-Systeme (SPS) auf Basis von Internet-of-Things (IoT)-Technologien und Cloud-Architekturen zunehmend an Bedeutung. Trotz der vielversprechenden Möglichkeiten zur Reduktion von Emissionen und Lärm standen Parkhausbetreiber vor der Herausforderung, die wirtschaftliche Rentabilität solcher Systeme verlässlich zu bewerten. Insbesondere die laufenden Kosten für Sensorik und Cloud-Dienste im Verhältnis zum erwarteten Mehrwert stellten eine schwer kalkulierbare Hürde für Investitionsentscheidungen dar. Um diese Problematik zu adressieren, wurde in dieser Arbeit der methodische Ansatz \enquote{Drive \& Decide} entwickelt und evaluiert. Dafür wurde ein prototypisches, IoT-basiertes SPS mit serverloser AWS-Architektur -- bestehend aus einem physischen Parkmodell und einer virtuellen Simulation -- entworfen, um darauf aufbauend ein detailliertes Kosten-Nutzen-Modell zur Erfassung von Investitionen (CAPEX) und nutzungsabhängigen Cloud-Kosten (OPEX) zu erstellen. Die Auswertung zeigte, dass durch diesen ereignisgetriebenen Aufbau belastbare betriebswirtschaftliche Kennzahlen wie die Total Cost of Ownership (TCO), der Return on Investment (ROI) und der Break-Even-Punkt präzise abgeleitet werden konnten. Der vorgeschlagene Ansatz lieferte somit einen messbaren Mehrwert in Form einer datengetriebenen und fundierten Entscheidungsgrundlage für Parkhausbetreiber. Im größeren Kontext verdeutlichten die Ergebnisse, dass eine flächendeckende Einführung von Smart-Parking-Lösungen bei korrekter Planung wirtschaftlich nachhaltig umsetzbar war. Damit wurde eine essenzielle Grundlage für effizientere urbane Mobilitätskonzepte geschaffen, welche zukünftig die Parksuche erleichtern und die Lebensqualität in Städten spürbar verbessern können.

\end{abstract}